%% un-sdg.tex: UN Sustainable Development Goals (FEUP regulations)
%% -------------------------------------------------------

\chapter*{UN Sustainable Development Goals} \label{chap:UnitedNations}

The United Nations Sustainable Development Goals (SDGs) provide a global framework to achieve a better and more sustainable future for all.
They include 17 goals addressing major challenges such as poverty, inequality, climate change, environmental degradation, peace, and justice.

This dissertation contributes primarily to \textbf{SDG 14 (Life Below Water)} by improving the efficiency and scientific value of in situ ocean observations through closed-loop, multi-agent adaptive sampling.
By enabling AUV teams to identify and sample \emph{dynamically distinct regimes} (e.g., fronts, water masses, plume-influenced areas) using machine learning, the work supports better monitoring and modelling of coastal environments.
A secondary contribution exists to \textbf{SDG 13 (Climate Action)}, since more informative and better-targeted ocean observations can improve downstream forecasting and understanding of climate-relevant ocean variability.

The specific Sustainable Development Goals mentioned have the following names:

\begin{description}
\item [SDG 13] Take urgent action to combat climate change and its impacts
\item [SDG 14] Conserve and sustainably use the oceans, seas and marine resources for sustainable development
\end{description}

\begin{center}
\begin{tabular}{|l|l|p{58mm}|p{52mm}|}
\hline
\textbf{SDG} & \textbf{Target} & \textbf{Contribution} & \textbf{Performance Indicators and Metrics} \\
\hline
\hline

\multirow{2}{*}{14} &
14.2 &
Enable more efficient monitoring of coastal ecosystem dynamics by directing multi-AUV sampling toward regime boundaries (fronts/plumes) and under-sampled regimes, improving the spatial/temporal coverage of observations for ecosystem assessment and management. &
(i) Coverage of high-variability/boundary regions (\% of boundary length/area sampled); (ii) reduction in out-of-regime prediction error (e.g., RMSE/MAE split by regime); (iii) increase in information gain or uncertainty reduction per kilometre/time. \\
\cline{2-4}

&
14.a &
Support marine scientific research capacity by providing an onboard-feasible closed-loop pipeline (regime inference + informative planning + coordination) and a benchmarking protocol for comparing methods under realistic communication and flow constraints. &
(i) Runtime/memory footprint of onboard components; (ii) reproducible benchmark suite released (scenarios, baselines, metrics); (iii) performance vs. baselines: redundancy overlap reduction (\%), utility per vehicle, and coordination robustness under dropouts/latency. \\
\hline

\multirow{2}{*}{13} &
13.1 &
Improve resilience of ocean observation strategies to heterogeneous and time-varying conditions by using regime-aware inference and update logic, reducing harmful cross-regime coupling in model updates and supporting more robust monitoring during rapid environmental changes. &
(i) Robustness under distribution shift: degradation in regime-classification confidence/calibration; (ii) performance under stress tests (currents, comm dropouts): mission success rate and constraint violations; (iii) in-regime vs out-of-regime forecast/update error change. \\
\cline{2-4}

&
13.3 &
Contribute to improved understanding and methodological practice for climate-relevant coastal variability by linking observations to uncertainty-aware decision-making and explicit evaluation of forecast-improvement proxies (when available). &
(i) Improvement in forecast-error proxy metrics after incorporating collected observations; (ii) uncertainty calibration metrics (e.g., reliability of confidence/entropy vs boundary regions); (iii) documented methodological guidelines for regime-aware sampling and evaluation. \\
\hline

\end{tabular}
\end{center}
